\chapter{Conclusões}
\label{cap:conclusoes}

\section{Considerações Finais}

Retomar e resumir, sucintamente, a motivação do trabalho, bem como a metodologia, os resultados obtidos e as contribuições alcançadas por este trabalho de TCC.

Discuta especificamente o \textbf{alcance dos objetivos} geral e específicos.

\section{Síntese da Experiência na Fábrica de Software}

Apresente uma síntese da experiência vivenciada, destacando os aspectos mais relevantes do processo, os principais resultados e a importância da participação no projeto para o desenvolvimento acadêmico e profissional do aluno. 

Discutir as principais lições aprendidas, sintetizando respostas para as seguintes perguntas:
\begin{enumerate}
    \item Quais são as suas impressões sobre a experiência realizada de prática de Engenharia de Software em um projeto real da Fábrica de Software?
    \item Quais os principais aprendizados, conquistas e evoluções em competências pessoais/comportamentais (\textit{soft skills}) e técnicas (\textit{hard skills}) obtidos da experiência realizada?
    \item Há outros benefícios associados à esta experiência, além de melhoria de competências \textit{soft} e \textit{hard}?
    \item Quais as principais dificuldades enfrentadas na experiência prática vivenciada?
    \item Quais sugestões você daria para estudantes que irão vivenciar esta experiência nos próximos semestres?
\end{enumerate}

\section{Trabalhos Relacionados}

Esta seção deve referenciar estudos, artigos, relatos de experiência e outros trabalhos acadêmicos que abordam temas semelhantes ao deste TCC. 

Deve ser feita uma análise comparativa entre os trabalhos referenciados e o presente documento de TCC. A comparação com trabalhos existentes permite situar o presente relato no contexto da literatura, identificar contribuições relevantes e confirmar ou refutar resultados e hipóteses apresentadas em outros trabalhos.
 
\section{Trabalhos Futuros}

Discuta as limitações do presente trabalho, indicando como futuros trabalhos podem reduzir essas limitações. Também aponte possíveis melhorias ou perspectivas para trabalhos futuros que complementem os resultados gerados neste TCC. Por exemplo, indique potenciais melhorias no produto de software que foi alvo dos processos de desenvolvimento ou manutenção, possíveis melhorias nesses processos, ou  ainda possíveis aplicações de ferramentas ou tecnologias que poderiam melhorar o produto ou o processo de Engenharia de Software.